% ===== PRÉAMBULE CANONIQUE — à coller VERBATIM en tête de CHAQUE .tex =====
\documentclass[11pt,a4paper]{article}
\usepackage[utf8]{inputenc}\usepackage[T1]{fontenc}\usepackage{lmodern}
\usepackage[french]{babel}
\usepackage{amsmath,amssymb,mathtools}\usepackage{icomma}
\usepackage[version=4]{mhchem}          % \ce{...} : équations & formules
\usepackage{siunitx}\sisetup{locale=FR} % \qty{}{}, \si{}, \ang{}
\usepackage{chemfig}                    % \chemfig{...} : structures développées
\usepackage{xcolor}\usepackage{colortbl}
\usepackage{tikz}\usepackage{pgfplots}\pgfplotsset{compat=1.18}
\usepackage[most]{tcolorbox}\usepackage{enumitem}\usepackage{array,booktabs}
\usepackage[a4paper,margin=2.2cm]{geometry}
\usetikzlibrary{arrows.meta,calc,angles,quotes,positioning,patterns,backgrounds,shapes.geometric}
\definecolor{primary}{RGB}{222,49,99}\definecolor{primarydark}{RGB}{150,22,55}
\definecolor{defcol}{RGB}{0,114,178}\definecolor{methcol}{RGB}{52,92,148}
\definecolor{excol}{RGB}{0,135,90}\definecolor{attncol}{RGB}{200,40,55}
\tcbset{boxbase/.style={enhanced,breakable,boxrule=0.9pt,arc=2.5pt,
  left=3mm,right=3mm,top=2.3mm,bottom=2.3mm,fonttitle=\bfseries,coltitle=white}}
% --- boîtes TUTORIEL (noms EXACTS, ne pas renommer) ---
\newtcolorbox{cle}[1][]{boxbase,colback=defcol!6,colframe=defcol,
  title={\textsc{À retenir}\ifx\relax#1\relax\relax\else\textmd{ : \itshape #1}\fi}}
\newtcolorbox{methode}[1][]{boxbase,colback=methcol!7,colframe=methcol,sharp corners,
  title={\textsc{Méthode}\ifx\relax#1\relax\relax\else\textmd{ --- \itshape #1}\fi}}
\newtcolorbox{exemplebox}[1][]{boxbase,colback=excol!5,colframe=excol,
  title={\textsc{Exemple}\ifx\relax#1\relax\relax\else\textmd{ : \itshape #1}\fi}}
\newtcolorbox{piege}[1][]{boxbase,colback=attncol!6,colframe=attncol,
  title={\textsc{Piège}\ifx\relax#1\relax\relax\else\textmd{ : \itshape #1}\fi}}
\newtcolorbox{remarque}[1][]{boxbase,colback=black!4,colframe=black!45,
  title={\textsc{Remarque}\ifx\relax#1\relax\relax\else\textmd{ : \itshape #1}\fi}}
% --- boîtes EXERCICE / CORRIGÉ (noms EXACTS, auto-comptées) ---
\newtcolorbox[auto counter]{exo}[1][]{boxbase,colback=primary!4,colframe=primary,
  title={\textsc{Exercice}~\thetcbcounter\ifx\relax#1\relax\relax\else\textmd{ --- \itshape #1}\fi}}
\newtcolorbox[auto counter]{corrige}[1][]{boxbase,colback=excol!4,colframe=excol,
  title={\textsc{Corrigé}~\thetcbcounter\ifx\relax#1\relax\relax\else\textmd{ --- \itshape #1}\fi}}
\newcommand{\R}{\mathbb{R}}\newcommand{\N}{\mathbb{N}}\newcommand{\Z}{\mathbb{Z}}\newcommand{\ds}{\displaystyle}
% (un \usepackage{fancyhdr}+en-tête est possible ; il est ignoré côté web)
% =========================================================================
\begin{document}

\begin{corrige}[l'échelle des oxacides du chlore]
Du n.o.\ le plus bas au plus haut : hypo-\ldots-eux, \ldots-eux, \ldots-ique, per-\ldots-ique.
\begin{enumerate}[label=\alph*)]
  \item \ce{HClO} : acide hypochloreux.
  \item \ce{HClO2} : acide chloreux.
  \item \ce{HClO3} : acide chlorique.
  \item \ce{HClO4} : acide perchlorique.
\end{enumerate}
\end{corrige}

\begin{corrige}[de l'oxacide à son anion : nom et formule]
On applique -eux$\to$-ite, -ique$\to$-ate ; les préfixes se gardent.
\begin{enumerate}[label=\alph*)]
  \item hypochlor\emph{eux} $\to$ hypochlor\emph{ite} \ce{ClO-}.
  \item chlor\emph{ique} $\to$ chlor\emph{ate} \ce{ClO3-}.
  \item perchlor\emph{ique} $\to$ perchlor\emph{ate} \ce{ClO4-}.
\end{enumerate}
\end{corrige}

\begin{corrige}[de l'anion à l'oxacide]
On remonte -ite$\to$-eux, -ate$\to$-ique, puis on ajoute autant de \ce{H} que la valence de l'anion.
\begin{enumerate}[label=\alph*)]
  \item nitr\emph{ite} \ce{NO2-} $\to$ acide nitr\emph{eux} \ce{HNO2}.
  \item sulf\emph{ate} \ce{SO4^2-} $\to$ acide sulfur\emph{ique} \ce{H2SO4}.
  \item carbon\emph{ate} \ce{CO3^2-} $\to$ acide carbon\emph{ique} \ce{H2CO3}.
\end{enumerate}
\end{corrige}

\begin{corrige}[transposer l'échelle au brome]
\begin{enumerate}[label=\alph*)]
  \item \ce{HBrO} : acide hypobromeux.
  \item \ce{HBrO3} : acide bromique.
  \item \ce{HBrO4} : acide perbromique.
\end{enumerate}
\end{corrige}

\begin{corrige}[distinguer -ite et -ate : écrire la formule de l'anion]
Le suffixe fixe le nombre d'oxygènes (l'\emph{ate} en a un de plus que l'\emph{ite}).
\begin{enumerate}[label=\alph*)]
  \item sulfite : \ce{SO3^2-}.
  \item sulfate : \ce{SO4^2-}.
  \item nitrite : \ce{NO2-}.
  \item nitrate : \ce{NO3-}.
\end{enumerate}
\end{corrige}

\end{document}
